\chapter{Background}

Demand shaping has long been a central challenge in operations, retail, and
energy systems. Traditional approaches rely on static incentives, broad
discounting, or heuristic rules that fail to account for individual behavioral
dynamics. This chapter reviews the theoretical foundations that motivate the
Behavioral Engine for Demand (BED), including behavioral economics, habit
formation, reinforcement learning, and constrained decision-making.

\section{Behavioral Economics and Incentive Response}

Behavioral economics provides a framework for understanding how individuals
respond to incentives in ways that deviate from classical rational models.
Empirical studies consistently show that consumers exhibit:

\begin{itemize}
    \item \textbf{Recency effects}: recent actions disproportionately influence
    future behavior.
    \item \textbf{Habit formation}: repeated actions increase the likelihood of
    repetition, independent of utility.
    \item \textbf{Responsiveness heterogeneity}: individuals differ in their
    sensitivity to incentives.
\end{itemize}

These effects motivate the BED state representation, which tracks recency,
habit, and responsiveness as latent behavioral variables.

\section{Demand Shaping in Retail and Energy Systems}

Demand shaping has been studied in domains such as:

\begin{itemize}
    \item \textbf{Retail}: loyalty programs, coupons, targeted promotions.
    \item \textbf{Energy}: time-of-use pricing, demand response programs.
    \item \textbf{Transportation}: congestion pricing, mobility incentives.
\end{itemize}

However, most systems rely on static or rule-based policies. They do not adapt
to individual behavioral trajectories, nor do they optimize incentives under
budget constraints. BED addresses these limitations by introducing a closed-loop
policy that learns from behavioral state transitions.

\section{Reinforcement Learning Foundations}

Reinforcement learning (RL) provides a mathematical framework for sequential
decision-making under uncertainty. In RL, an agent observes a state, selects an
action, and receives a reward. BED draws from RL in three ways:

\begin{enumerate}
    \item \textbf{State representation}: recency, habit, and responsiveness form
    a Markovian behavioral state.
    \item \textbf{Transition dynamics}: incentives influence state evolution.
    \item \textbf{Policy optimization}: actions are chosen to maximize long-term
    behavioral outcomes.
\end{enumerate}

Unlike classical RL, BED must operate under strict budget constraints and
fairness requirements, motivating the use of constrained MDPs.

\section{Constrained Markov Decision Processes}

A Constrained Markov Decision Process (CMDP) extends the MDP framework by
introducing resource constraints. In BED, the primary constraint is the
incentive budget. The CMDP formulation ensures that:

\begin{itemize}
    \item incentives are allocated efficiently,
    \item long-term behavioral impact is maximized,
    \item fairness and exposure constraints are respected.
\end{itemize}

This structure allows BED to balance short-term engagement with long-term
behavioral shaping.

\section{Summary}

The BED framework integrates insights from behavioral economics, demand shaping,
reinforcement learning, and constrained optimization. These foundations
establish the need for a system that models individual behavior, updates
dynamically, and allocates incentives optimally under constraints. The next
chapter introduces the system architecture that operationalizes these ideas.
